% --- Archon physics formalization source begin ---
% archon:physics
% archon:covers PhyXMiniProblems/problem_phyx_mini_0942.lean
% archon:source-report reports/phyx_mini/problem_phyx_mini_0942.source.json

\chapter{Physics problem phyx\_mini\_0942}
\label{ch:PhyXMiniProblems_problem_phyx_mini_0942}

\paragraph{Problem source.}
\#\# Physical scenario

A 200\textbackslash{} \textbackslash{}Omega\textbackslash{} resistor is connected in series with a 5.0\textbackslash{} \textbackslash{}mu\textbackslash{}mathrm\{F\} capacitor. The voltage across the resistor is \textbackslash{}( v\_R = (1.20\textbackslash{} \textbackslash{}mathrm\{V\}) \textbackslash{}cos[(2500\textbackslash{} \textbackslash{}mathrm\{rad/s\})t] \textbackslash{}).

\#\# Question

Determine the capacitive reactance of the capacitor.

\#\# Answer choices

- A: A: 20\textbackslash{} \textbackslash{}Omega\textbackslash{}

- B: B: 80\textbackslash{} \textbackslash{}Omega\textbackslash{}

- C: C: 100\textbackslash{} \textbackslash{}Omega\textbackslash{}

- D: D: 200\textbackslash{} \textbackslash{}Omega\textbackslash{}

\#\# Figure caption (auxiliary; use the image as primary evidence)

The image is a schematic diagram of an electrical circuit. Here are the components and details:

1. **AC Voltage Source**: Represented by a circle with a sine wave inside, indicating an alternating current source.

2. **Components**:
   - **Capacitor (C)**: Labeled with a capacitance of 5.0 µF (microfarads).
   - **Resistor (R)**: Labeled with a resistance of 200 Ω (ohms).

3. **Circuit Paths**:
   - Current \textbackslash{}( i \textbackslash{}) is indicated by an arrow pointing to the left.
   - Voltage across the capacitor is denoted as \textbackslash{}( v\_c \textbackslash{}).
   - Voltage across the resistor is denoted as \textbackslash{}( v\_R \textbackslash{}).

4. **Expression for Voltage**:
   - \textbackslash{}( v\_R = (1.20 \textbackslash{}text\{ V\}) \textbackslash{}cos(2500 \textbackslash{}text\{ rad/s\})t \textbackslash{})
   - This represents the voltage across the resistor as a function of time with a cosine waveform, an amplitude of 1.20 volts, and an angular frequency of 2500 radians per second.

5. **Connections**:
   - Wires connect the voltage source, capacitor, and resistor in series, forming a closed loop.

This diagram illustrates the relationships and components in an AC circuit with given electrical properties.

\#\# Dataset metadata

Electric Circuits; Physical Model Grounding Reasoning, Numerical Reasoning

\paragraph{Recorded answer/context.}
B: B: 80\textbackslash{} \textbackslash{}Omega\textbackslash{}

\paragraph{Figure/image path.}
/root/proposal\_for\_physic/science-mango/phyx\_mini\_run/phyx\_data/test\_image/942.png

\paragraph{Formalization target.}
create a compiling Lean file with sorry bodies at `PhyXMiniProblems/problem\_phyx\_mini\_0942.lean`. The Lean declarations must preserve the physical quantities, dimensions or dimensional roles, figure labels, governing-law hypotheses, and final relation expressed by this problem.
Use Mathlib/Physlib names found through LeanExplore where available. If a physics API is missing, introduce faithful local abstractions rather than scalar placeholder aliases.

\begin{theorem}[Physics formalization target]
\label{thm:physics:phyx_mini_0942:target}
\lean{PhyXMiniProblems.ProblemPhyXMini0942.capacitive_reactance_eq_eighty_ohms}
\uses{def:physics:phyx-mini-0942:phyxminiproblems-problemphyxmini0942-resistanceinohms, def:physics:phyx-mini-0942:phyxminiproblems-problemphyxmini0942-seriesrccircuitsetup, def:physics:phyx-mini-0942:phyxminiproblems-problemphyxmini0942-matchesseriesrcscenario, def:physics:phyx-mini-0942:phyxminiproblems-problemphyxmini0942-matchesprimaryseriesrcfigure, def:physics:phyx-mini-0942:phyxminiproblems-problemphyxmini0942-hasphysicalseriesrcparameters, def:physics:phyx-mini-0942:phyxminiproblems-problemphyxmini0942-satisfiesidealcapacitivereactancelaw}
The capacitor's reactance magnitude is $80$ ohms.
\end{theorem}
\begin{proof}
The figure gives $C=5$ microfarads and angular frequency $\omega=2500$ radians per second.  Convert the capacitance to $5/10^6$ farads and substitute both readouts into $X_C=1/(\omega C)$.  The denominator is $2500\cdot5/10^6=1/80$, hence $X_C=80$ ohms.
\end{proof}

% --- Archon declaration topology begin ---
\section{Declaration topology}
\begin{definition}[electric Potential Dimension]
  \label{def:physics:phyx-mini-0942:phyxminiproblems-problemphyxmini0942-electricpotentialdimension}
  \lean{PhyXMiniProblems.ProblemPhyXMini0942.electricPotentialDimension}
  The dimension M L² T⁻² C⁻¹ of electric potential (volt)
\end{definition}
\begin{proof}
  This is the declared model definition of electric Potential Dimension.
\end{proof}
\begin{definition}[electric Current Dimension]
  \label{def:physics:phyx-mini-0942:phyxminiproblems-problemphyxmini0942-electriccurrentdimension}
  \lean{PhyXMiniProblems.ProblemPhyXMini0942.electricCurrentDimension}
  The dimension C T⁻¹ of electric current (ampere)
\end{definition}
\begin{proof}
  This is the declared model definition of electric Current Dimension.
\end{proof}
\begin{definition}[electrical Resistance Dimension]
  \label{def:physics:phyx-mini-0942:phyxminiproblems-problemphyxmini0942-electricalresistancedimension}
  \lean{PhyXMiniProblems.ProblemPhyXMini0942.electricalResistanceDimension}
  The dimension M L² T⁻¹ C⁻² of electrical resistance (ohm)
\end{definition}
\begin{proof}
  This is the declared model definition of electrical Resistance Dimension.
\end{proof}
\begin{definition}[capacitance Dimension]
  \label{def:physics:phyx-mini-0942:phyxminiproblems-problemphyxmini0942-capacitancedimension}
  \lean{PhyXMiniProblems.ProblemPhyXMini0942.capacitanceDimension}
  The dimension M⁻¹ L⁻² T² C² of capacitance (farad)
\end{definition}
\begin{proof}
  This is the declared model definition of capacitance Dimension.
\end{proof}
\begin{definition}[Resistance Magnitude]
  \label{def:physics:phyx-mini-0942:phyxminiproblems-problemphyxmini0942-resistancemagnitude}
  \lean{PhyXMiniProblems.ProblemPhyXMini0942.ResistanceMagnitude}
  \uses{def:physics:phyx-mini-0942:phyxminiproblems-problemphyxmini0942-electricalresistancedimension}
  A nonnegative, unit-independent electrical resistance magnitude
\end{definition}
\begin{proof}
  This is the declared model definition of Resistance Magnitude.
\end{proof}
\begin{definition}[Capacitance Magnitude]
  \label{def:physics:phyx-mini-0942:phyxminiproblems-problemphyxmini0942-capacitancemagnitude}
  \lean{PhyXMiniProblems.ProblemPhyXMini0942.CapacitanceMagnitude}
  \uses{def:physics:phyx-mini-0942:phyxminiproblems-problemphyxmini0942-capacitancedimension}
  A nonnegative, unit-independent capacitance magnitude
\end{definition}
\begin{proof}
  This is the declared model definition of Capacitance Magnitude.
\end{proof}
\begin{definition}[Angular Frequency Magnitude]
  \label{def:physics:phyx-mini-0942:phyxminiproblems-problemphyxmini0942-angularfrequencymagnitude}
  \lean{PhyXMiniProblems.ProblemPhyXMini0942.AngularFrequencyMagnitude}
  A nonnegative, unit-independent angular-frequency magnitude
\end{definition}
\begin{proof}
  This is the declared model definition of Angular Frequency Magnitude.
\end{proof}
\begin{definition}[Electric Potential Quantity]
  \label{def:physics:phyx-mini-0942:phyxminiproblems-problemphyxmini0942-electricpotentialquantity}
  \lean{PhyXMiniProblems.ProblemPhyXMini0942.ElectricPotentialQuantity}
  \uses{def:physics:phyx-mini-0942:phyxminiproblems-problemphyxmini0942-electricpotentialdimension}
  A signed, unit-independent electric-potential quantity
\end{definition}
\begin{proof}
  This is the declared model definition of Electric Potential Quantity.
\end{proof}
\begin{definition}[Electric Current Quantity]
  \label{def:physics:phyx-mini-0942:phyxminiproblems-problemphyxmini0942-electriccurrentquantity}
  \lean{PhyXMiniProblems.ProblemPhyXMini0942.ElectricCurrentQuantity}
  \uses{def:physics:phyx-mini-0942:phyxminiproblems-problemphyxmini0942-electriccurrentdimension}
  A signed, unit-independent electric-current quantity
\end{definition}
\begin{proof}
  This is the declared model definition of Electric Current Quantity.
\end{proof}
\begin{definition}[Time Quantity]
  \label{def:physics:phyx-mini-0942:phyxminiproblems-problemphyxmini0942-timequantity}
  \lean{PhyXMiniProblems.ProblemPhyXMini0942.TimeQuantity}
  A signed, unit-independent time coordinate
\end{definition}
\begin{proof}
  This is the declared model definition of Time Quantity.
\end{proof}
\begin{definition}[nonnegative SIReadout]
  \label{def:physics:phyx-mini-0942:phyxminiproblems-problemphyxmini0942-nonnegativesireadout}
  \lean{PhyXMiniProblems.ProblemPhyXMini0942.nonnegativeSIReadout}
  Coherent-SI readout of a nonnegative dimensionful quantity
\end{definition}
\begin{proof}
  This is the declared model definition of nonnegative SIReadout.
\end{proof}
\begin{definition}[signed SIReadout]
  \label{def:physics:phyx-mini-0942:phyxminiproblems-problemphyxmini0942-signedsireadout}
  \lean{PhyXMiniProblems.ProblemPhyXMini0942.signedSIReadout}
  Coherent-SI readout of a signed dimensionful quantity
\end{definition}
\begin{proof}
  This is the declared model definition of signed SIReadout.
\end{proof}
\begin{definition}[resistance In Ohms]
  \label{def:physics:phyx-mini-0942:phyxminiproblems-problemphyxmini0942-resistanceinohms}
  \lean{PhyXMiniProblems.ProblemPhyXMini0942.resistanceInOhms}
  \uses{def:physics:phyx-mini-0942:phyxminiproblems-problemphyxmini0942-resistancemagnitude, def:physics:phyx-mini-0942:phyxminiproblems-problemphyxmini0942-nonnegativesireadout}
  Coherent-SI resistance readout, in ohms
\end{definition}
\begin{proof}
  This is the declared model definition of resistance In Ohms.
\end{proof}
\begin{definition}[capacitance In Farads]
  \label{def:physics:phyx-mini-0942:phyxminiproblems-problemphyxmini0942-capacitanceinfarads}
  \lean{PhyXMiniProblems.ProblemPhyXMini0942.capacitanceInFarads}
  \uses{def:physics:phyx-mini-0942:phyxminiproblems-problemphyxmini0942-capacitancemagnitude, def:physics:phyx-mini-0942:phyxminiproblems-problemphyxmini0942-nonnegativesireadout}
  Coherent-SI capacitance readout, in farads
\end{definition}
\begin{proof}
  This is the declared model definition of capacitance In Farads.
\end{proof}
\begin{definition}[angular Frequency In Radians Per Second]
  \label{def:physics:phyx-mini-0942:phyxminiproblems-problemphyxmini0942-angularfrequencyinradianspersecond}
  \lean{PhyXMiniProblems.ProblemPhyXMini0942.angularFrequencyInRadiansPerSecond}
  \uses{def:physics:phyx-mini-0942:phyxminiproblems-problemphyxmini0942-angularfrequencymagnitude, def:physics:phyx-mini-0942:phyxminiproblems-problemphyxmini0942-nonnegativesireadout}
  Coherent-SI angular-frequency readout, in radians per second
\end{definition}
\begin{proof}
  This is the declared model definition of angular Frequency In Radians Per Second.
\end{proof}
\begin{definition}[electric Potential In Volts]
  \label{def:physics:phyx-mini-0942:phyxminiproblems-problemphyxmini0942-electricpotentialinvolts}
  \lean{PhyXMiniProblems.ProblemPhyXMini0942.electricPotentialInVolts}
  \uses{def:physics:phyx-mini-0942:phyxminiproblems-problemphyxmini0942-electricpotentialquantity, def:physics:phyx-mini-0942:phyxminiproblems-problemphyxmini0942-signedsireadout}
  Coherent-SI electric-potential readout, in volts
\end{definition}
\begin{proof}
  This is the declared model definition of electric Potential In Volts.
\end{proof}
\begin{definition}[electric Current In Amperes]
  \label{def:physics:phyx-mini-0942:phyxminiproblems-problemphyxmini0942-electriccurrentinamperes}
  \lean{PhyXMiniProblems.ProblemPhyXMini0942.electricCurrentInAmperes}
  \uses{def:physics:phyx-mini-0942:phyxminiproblems-problemphyxmini0942-electriccurrentquantity, def:physics:phyx-mini-0942:phyxminiproblems-problemphyxmini0942-signedsireadout}
  Coherent-SI electric-current readout, in amperes
\end{definition}
\begin{proof}
  This is the declared model definition of electric Current In Amperes.
\end{proof}
\begin{definition}[time In Seconds]
  \label{def:physics:phyx-mini-0942:phyxminiproblems-problemphyxmini0942-timeinseconds}
  \lean{PhyXMiniProblems.ProblemPhyXMini0942.timeInSeconds}
  \uses{def:physics:phyx-mini-0942:phyxminiproblems-problemphyxmini0942-timequantity, def:physics:phyx-mini-0942:phyxminiproblems-problemphyxmini0942-signedsireadout}
  Coherent-SI time-coordinate readout, in seconds
\end{definition}
\begin{proof}
  This is the declared model definition of time In Seconds.
\end{proof}
\begin{definition}[Circuit Component]
  \label{def:physics:phyx-mini-0942:phyxminiproblems-problemphyxmini0942-circuitcomponent}
  \lean{PhyXMiniProblems.ProblemPhyXMini0942.CircuitComponent}
  The three components visible in the supplied circuit schematic
\end{definition}
\begin{proof}
  This is the declared model definition of Circuit Component.
\end{proof}
\begin{definition}[Voltage Symbol]
  \label{def:physics:phyx-mini-0942:phyxminiproblems-problemphyxmini0942-voltagesymbol}
  \lean{PhyXMiniProblems.ProblemPhyXMini0942.VoltageSymbol}
  The two explicitly named component-voltage symbols in the figure
\end{definition}
\begin{proof}
  This is the declared model definition of Voltage Symbol.
\end{proof}
\begin{definition}[Horizontal Direction]
  \label{def:physics:phyx-mini-0942:phyxminiproblems-problemphyxmini0942-horizontaldirection}
  \lean{PhyXMiniProblems.ProblemPhyXMini0942.HorizontalDirection}
  Horizontal directions used by the current arrow in the schematic
\end{definition}
\begin{proof}
  This is the declared model definition of Horizontal Direction.
\end{proof}
\begin{definition}[Series RCFigure]
  \label{def:physics:phyx-mini-0942:phyxminiproblems-problemphyxmini0942-seriesrcfigure}
  \lean{PhyXMiniProblems.ProblemPhyXMini0942.SeriesRCFigure}
  \uses{def:physics:phyx-mini-0942:phyxminiproblems-problemphyxmini0942-circuitcomponent, def:physics:phyx-mini-0942:phyxminiproblems-problemphyxmini0942-voltagesymbol, def:physics:phyx-mini-0942:phyxminiproblems-problemphyxmini0942-horizontaldirection}
  Literal information visible in image 942.png. Numerical fields record the numbers printed beside the physical quantities; their units are specified by their field names and tied to dimensionful quantities below
\end{definition}
\begin{proof}
  This is the declared model definition of Series RCFigure.
\end{proof}
\begin{definition}[Series RCCircuit Setup]
  \label{def:physics:phyx-mini-0942:phyxminiproblems-problemphyxmini0942-seriesrccircuitsetup}
  \lean{PhyXMiniProblems.ProblemPhyXMini0942.SeriesRCCircuitSetup}
  \uses{def:physics:phyx-mini-0942:phyxminiproblems-problemphyxmini0942-resistancemagnitude, def:physics:phyx-mini-0942:phyxminiproblems-problemphyxmini0942-capacitancemagnitude, def:physics:phyx-mini-0942:phyxminiproblems-problemphyxmini0942-angularfrequencymagnitude, def:physics:phyx-mini-0942:phyxminiproblems-problemphyxmini0942-electricpotentialquantity, def:physics:phyx-mini-0942:phyxminiproblems-problemphyxmini0942-electriccurrentquantity, def:physics:phyx-mini-0942:phyxminiproblems-problemphyxmini0942-timequantity, def:physics:phyx-mini-0942:phyxminiproblems-problemphyxmini0942-seriesrcfigure}
  The physical objects and time-dependent observables in the displayed circuit. The capacitive reactance is independent data here, rather than an expansion of the answer formula
\end{definition}
\begin{proof}
  This is the declared model definition of Series RCCircuit Setup.
\end{proof}
\begin{definition}[Matches Series RCScenario]
  \label{def:physics:phyx-mini-0942:phyxminiproblems-problemphyxmini0942-matchesseriesrcscenario}
  \lean{PhyXMiniProblems.ProblemPhyXMini0942.MatchesSeriesRCScenario}
  \uses{def:physics:phyx-mini-0942:phyxminiproblems-problemphyxmini0942-seriesrccircuitsetup}
  The prose and schematic specify a closed series circuit with an AC source
\end{definition}
\begin{proof}
  This is the declared model definition of Matches Series RCScenario.
\end{proof}
\begin{definition}[Matches Primary Series RCFigure]
  \label{def:physics:phyx-mini-0942:phyxminiproblems-problemphyxmini0942-matchesprimaryseriesrcfigure}
  \lean{PhyXMiniProblems.ProblemPhyXMini0942.MatchesPrimarySeriesRCFigure}
  \uses{def:physics:phyx-mini-0942:phyxminiproblems-problemphyxmini0942-resistanceinohms, def:physics:phyx-mini-0942:phyxminiproblems-problemphyxmini0942-capacitanceinfarads, def:physics:phyx-mini-0942:phyxminiproblems-problemphyxmini0942-angularfrequencyinradianspersecond, def:physics:phyx-mini-0942:phyxminiproblems-problemphyxmini0942-electricpotentialinvolts, def:physics:phyx-mini-0942:phyxminiproblems-problemphyxmini0942-timeinseconds, def:physics:phyx-mini-0942:phyxminiproblems-problemphyxmini0942-seriesrccircuitsetup}
  Primary-raster evidence: component labels, the current arrow, and every numerical annotation printed in the image. The waveform equality is a calibrated transcription of the displayed v\_R(t) expression. It does not state the requested reactance
\end{definition}
\begin{proof}
  This is the declared model definition of Matches Primary Series RCFigure.
\end{proof}
\begin{definition}[Has Physical Series RCParameters]
  \label{def:physics:phyx-mini-0942:phyxminiproblems-problemphyxmini0942-hasphysicalseriesrcparameters}
  \lean{PhyXMiniProblems.ProblemPhyXMini0942.HasPhysicalSeriesRCParameters}
  \uses{def:physics:phyx-mini-0942:phyxminiproblems-problemphyxmini0942-resistanceinohms, def:physics:phyx-mini-0942:phyxminiproblems-problemphyxmini0942-capacitanceinfarads, def:physics:phyx-mini-0942:phyxminiproblems-problemphyxmini0942-angularfrequencyinradianspersecond, def:physics:phyx-mini-0942:phyxminiproblems-problemphyxmini0942-electricpotentialinvolts, def:physics:phyx-mini-0942:phyxminiproblems-problemphyxmini0942-seriesrccircuitsetup}
  Positivity conditions for the nondegenerate physical circuit data
\end{definition}
\begin{proof}
  This is the declared model definition of Has Physical Series RCParameters.
\end{proof}
\begin{definition}[Satisfies Ideal Capacitive Reactance Law]
  \label{def:physics:phyx-mini-0942:phyxminiproblems-problemphyxmini0942-satisfiesidealcapacitivereactancelaw}
  \lean{PhyXMiniProblems.ProblemPhyXMini0942.SatisfiesIdealCapacitiveReactanceLaw}
  \uses{def:physics:phyx-mini-0942:phyxminiproblems-problemphyxmini0942-resistanceinohms, def:physics:phyx-mini-0942:phyxminiproblems-problemphyxmini0942-capacitanceinfarads, def:physics:phyx-mini-0942:phyxminiproblems-problemphyxmini0942-angularfrequencyinradianspersecond, def:physics:phyx-mini-0942:phyxminiproblems-problemphyxmini0942-seriesrccircuitsetup}
  For an ideal capacitor driven sinusoidally at angular frequency ω, the capacitive-reactance magnitude is 1 / (ω C). This is the governing physical law, relating three independent observables; it does not assert the target numeric value
\end{definition}
\begin{proof}
  This is the declared model definition of Satisfies Ideal Capacitive Reactance Law.
\end{proof}
% --- Archon declaration topology end ---
% --- Archon physics formalization source end ---
